\section{Simulated Vehicle and Transferable Navigation}
\label{sec:vehicle}

\subsection{BlueROV2 Heavy in Simulation}

The simulated vehicle is the simulator's native BlueROV2 agent, which is the
Heavy configuration: eight thrusters in a vectored six-degree-of-freedom
layout matching the physical vehicle's frame configuration, with a
neutrally-buoyant 11.5\,kg rigid body, a centre of buoyancy 5\,cm above the
centre of mass, and a per-thruster force limit of $\pm 28.75$\,N. Motion is
produced by engine rigid-body physics under commanded thruster forces; the
kinematic teleport used by the pre-scientific version of this stack is
excluded from all scientific scenarios, which restores physical contact as an
observable event. Because the physical vehicle's frame configuration is the
same vectored layout, the actuation topology is matched by construction
rather than by tuning.

One implementation detail is worth recording because it invalidated an
earlier controller. The thruster geometry documented in the simulator's
Python agent interface is inconsistent with the engine implementation, and
using it produced inverted surge, absent heave, and spurious yaw. Open-loop
force probes together with the engine source established the correct
client-frame geometry: four vertical thrusters at $(\pm 0.12, \pm 0.2181,
+0.0809)$\,m and four horizontal $45^{\circ}$ pods at $(\pm 0.1562, \pm
0.0988, 0)$\,m, in a body frame with $x$ forward, $y$ left, $z$ up. The
corrected geometry reproduces every probe observation.

\subsection{Command Path and Low-Level Control}

The planner's output is a body-velocity command, and everything between that
command and the water is deliberately modeled as a chain of separately
calibratable stages (Fig.~\ref{fig:control}). A commanded twist passes
through a configurable command-latency buffer, a setpoint rate limiter, a
proportional--integral body-velocity controller for surge, sway, and heave, a
proportional--integral yaw-rate loop, and a proportional--derivative
roll/pitch stabilizer; the resulting body wrench is mapped to eight thruster
forces by a pseudo-inverse allocation with direction-preserving saturation. A
server-side watchdog zeroes the commanded velocity when no fresh command
arrives within 1\,s, mirroring the vehicle-side failsafe. The latency buffer
is the mechanism by which calibration level \STWO{} injects measured
command-path delay, and the controller gains and rate limits are what
\STHREE{} adjusts.

%==============================================================================
% Figure 2 --- BlueROV2 simulation/control architecture.  Status: READY.
% Source: docs/BLUEROV2_SIMULATION_DYNAMICS.md (architecture block),
% docs/BLUEROV2_HOLOOCEAN_INTEGRATION.md (thruster geometry, limits).
%==============================================================================
\begin{figure}[t]
\centering
\begin{tikzpicture}[
  font=\scriptsize,
  b/.style={draw, rounded corners=1pt, align=center, minimum height=5.4mm,
            text width=21mm, inner xsep=1.5pt},
  cal/.style={b, fill=yellow!20},
  ar/.style={-{Latex[length=1.5mm]}, thick},
  node distance=2.1mm,
]
\node[b] (cmd) {Planner twist\\$(u,v,w,r)_{\mathrm{cmd}}$};
\node[cal, below=of cmd] (lat) {Command latency\\buffer \textbf{(\STWO)}};
\node[cal, below=of lat] (rl)  {Setpoint rate\\limiter \textbf{(\STHREE)}};
\node[cal, below=of rl]  (pi)  {PI body-velocity\\+ PI yaw rate\\+ PD roll/pitch \textbf{(\STHREE)}};
\node[b, below=of pi]    (al)  {Pseudo-inverse\\8-thruster alloc.};
\node[b, below=of al]    (sat) {Direction-preserving\\saturation $\pm 28.75$\,N};
\node[b, below=of sat]   (ph)  {Rigid-body physics\\(BlueROV2 Heavy)};

\foreach \a/\b in {cmd/lat, lat/rl, rl/pi, pi/al, al/sat, sat/ph}
  \draw[ar] (\a) -- (\b);

% feedback
\draw[ar] (ph.east) -- ++(4mm,0) |- node[pos=0.25, right, font=\tiny,
  text width=16mm, align=left]
  {body velocity,\\rates, roll/pitch\\(simulated plant\\feedback only)} (pi.east);

% watchdog
\node[b, left=6mm of pi, text width=15mm, fill=red!8] (wd)
  {Watchdog\\(1\,s) $\rightarrow 0$};
\draw[ar] (wd) -- (pi);

\end{tikzpicture}
\caption{Simulated command path. Shaded stages are the ones the calibration
ladder adjusts: transport and processing delay at \STWO{}, and vehicle
response (rate limits, controller gains) at \STHREE{}. The velocity feedback
closing the inner loop belongs to the simulated \emph{plant}, standing in for
the physical vehicle's own thruster-loop behavior; it is not available to the
navigation estimator.}
\label{fig:control}
\end{figure}


Step responses of this simulated chain are reported in
Table~\ref{tab:stepresponse}. We stress what these numbers are and are not:
they characterize the \emph{simulated} vehicle and its controller at the
uncalibrated level \SZERO{}, and they are the reference against which the
measured physical responses will later be compared. They are not a
simulation-versus-reality result, and no such comparison exists yet.

%==============================================================================
% S0 simulated step responses.  Source: docs/BLUEROV2_SIMULATION_DYNAMICS.md,
% experiments/simulation/step_response_S0{,_run2,_run3}/manifest.json
% (3 fresh-process runs, bit-identical summaries).  Status: READY.
% NOTE: simulation-internal controller validation.  NOT a sim-vs-real result.
%==============================================================================
\begin{table}[t]
\centering
\caption{Step responses of the \emph{simulated} BlueROV2 Heavy at the
uncalibrated level \SZERO{}, from three consecutive fresh-process runs with
bit-identical summaries. These validate the simulated command path; they are
\emph{not} a simulation-versus-reality comparison, which requires the
measured physical responses acquired at \STHREE{}.}
\label{tab:stepresponse}
\footnotesize
\setlength{\tabcolsep}{4pt}
\renewcommand{\arraystretch}{1.15}
\begin{tabular}{@{}l r r r r@{}}
\toprule
\textbf{Axis} & \textbf{Amplitude} & \textbf{Rise 10--90\%} &
\textbf{SS error} & \textbf{Overshoot} \\
\midrule
Surge    & 0.40\,m/s   & 3.03\,s & 4.8\,\% & none \\
Sway     & 0.25\,m/s   & 2.90\,s & 4.5\,\% & none \\
Heave    & $-0.20$\,m/s& 2.03\,s & 3.2\,\% & none \\
Yaw rate & 0.30\,rad/s & 0.20\,s & 0.6\,\% & none \\
\bottomrule
\end{tabular}
\end{table}


\subsection{Transferable Navigation Without a Velocity Sensor}
\label{sec:navigation}

Because the physical vehicle has no Doppler velocity log, runtime navigation
must be reconstructed from signals that exist on both sides, and it is this
constraint --- not a modeling preference --- that fixes the estimator's form.
The runtime odometry propagates the \emph{commanded} body velocity through a
first-order vehicle-response model, with time constants $\tau_u = 1.2$\,s,
$\tau_v = 1.15$\,s, and $\tau_w = 0.8$\,s taken from the step-response
identification, and integrates it using the \emph{measured} yaw; the vertical
channel is the measured pressure depth rather than an integrated quantity.
When commands go stale beyond 1\,s the modeled velocity decays to zero,
mirroring the vehicle watchdog. The estimator is therefore a dead-reckoning
filter whose inputs are available identically in simulation and in water.

The corresponding restriction is enforced in code: the simulator's velocity
sensor feeds only the simulated low-level controller, where it plays the role
of the physical thruster loop's internal feedback, and never reaches the
navigation estimator, which refuses ground-truth-like topics at
construction. The cost of this design is drift, which we measure rather than
assume: over the accepted integration campaign the maximum dead-reckoning
error was 0.117\,m on straight runs and $\leq 0.20$\,m during avoidance
maneuvers (Sec.~\ref{sec:results}). The benefit is that the navigation
solution transfers without modification, so a sim-to-real difference in
avoidance outcome cannot be explained away by a navigation input that exists
only in simulation.
